% RESS Sections 3-4: setting and the four results, stated in words; proofs in Appendix A (proofs_body.tex).
\section{Demonstration setting}\label{sec:demo}

\paragraph{Unit, demand and failure} The unit under test is a fixed camera $s$ with a fixed detector. The camera observes
\emph{events}: maximal runs of frames that contain the target object. It samples frames with period $K$ and an unknown
phase $\varphi$, and runs the detector on the sampled frames. An event is \emph{missed} ($M=1$) if no sampled frame is a
hit; with $H$ the hit frames of an event and $\varphi$ uniform, $\PP(M=1)=1-|H\bmod K|/K$. The miss probability of camera
$s$ is $p_s$, and the demonstration goal is the assurance claim $p_s\le\varepsilon$ at confidence $1-\delta$.
$\UCP(x,n)$ denotes the one-sided $(1-\delta)$ Clopper--Pearson upper bound for $x$ failures in $n$ demands. We call an accepted assurance claim of this form a \emph{certificate}. In the language of acceptance sampling,
$\delta$ is the consumer's risk; the producer's risk is reported as the power of a plan at the observed miss rate.
All claims concern the miss probability averaged over the uniform sampling phase; a camera deployed with one fixed,
unknown phase is covered only on average over phases.

\paragraph{Simulated evidence and twins} A simulator produces data $Y\sim Q$. A \emph{paired twin} renders a synthetic
counterpart of each real event and runs the same detector at the same sampling phase, giving $M'$. We write
$q=\PP(M'=1)$ and $D=\PP(M=1,M'=0)$ (real miss, twin catch), the \emph{discordance}. Every rate in this paper is per
event; with three sprite variants per event, the twin catches an event if at least two variants do, and an event counts
as discordant if it is discordant at one or more sampling phases (worst phase; an integer count, conservative with
respect to the phase-averaged rate, which we give in parentheses). Table~\ref{tab:notation} collects the notation.

\paragraph{Random model} Scenes are drawn from a population; given a scene, events are drawn i.i.d.\ from the scene's
event law and each receives an independent uniform phase; given the event and the phase, the real and twin outcomes are
determined (fixed detector, fixed sprite seeds). Every per-event indicator below---real miss, twin miss, discordance at a
random phase, and the worst-phase indicator---is therefore a Bernoulli variable, independent across events given the
scene. In the audit, the ``truth'' of a scene is its empirical event law, resampled with replacement.

\begin{table}[t]
\centering
\caption{Notation.}
\label{tab:notation}
\small
\begin{tabular}{@{}L{0.26\textwidth}L{0.68\textwidth}@{}}
\toprule
symbol & meaning \\
\midrule
$\varepsilon,\ \delta$ & target miss probability; allowed probability of a false acceptance \\
$n_0(\varepsilon,\delta)$ & zero-failure sample size $\lceil\ln\delta/\ln(1-\varepsilon)\rceil$ \\
$K,\ \varphi,\ d_{\min}$ & sampling period, sampling phase, minimum event length (frames) \\
$M,\ M'$ & real miss, twin miss of an event at the same phase \\
$p_s,\ q_s,\ D_s$ & real miss probability, twin miss probability, discordance $\PP(M=1,M'=0)$ of camera $s$ \\
$m,\ n_s,\ N$ & calibration scenes, paired events of scene $s$, $N=\sum_sn_s$ \\
$X,\ J$ & discordant events; scenes with at least one discordant event \\
$\UCP(x,n)$ & one-sided $(1-\delta)$ Clopper--Pearson upper bound \\
fleet bound, population bound & $\UCP(X,N)$ and $\UCP(J,m)$ (Result~4) \\
\OPs, OP-A, OP-C & operating points $(\varepsilon,d_{\min},K)=(20\%,32,16)$, $(20\%,32,8)$, $(10\%,32,1)$ \\
\bottomrule
\end{tabular}
\end{table}

% =====================================================================================================
\section{What simulated evidence can buy: four results}\label{sec:results-theory}

The four results below are stated in words; \ref{app:proofs} gives the formal statements, the assumptions and the
proofs, each with a numerical check.

\paragraph{Result 1 (no free credit)}\label{res:c1} Suppose the simulated data are independent of the real demands and no
assumption links the simulator to reality, so that for every simulator law $Q$ both a perfect unit ($p=0$) and a unit just
worse than the target ($p\downarrow\varepsilon$) are possible. Then any acceptance procedure that is valid at confidence
$1-\delta$ accepts a perfect unit with probability at most $\delta/(1-\varepsilon)^n$ when it uses $n$ real demands,
whatever the simulator and however much simulation it uses. Accepting a perfect unit with certainty therefore still
needs $n\ge n_0(\varepsilon,\delta)$ real demands. The reason is that the simulated data have the same law under both
hypotheses, while an all-success real sample is possible under both. This makes precise, for the success-run test, the principle that external evidence cannot shorten a test whose validity must hold for every external law~\cite{koppschneider2020}, and the related limit on combining evidence for ultra-high dependability~\cite{littlewood1993}. A paired twin is built from the real events and is therefore \emph{not} an
independent simulator; pairing is what Result~3 exploits.

\paragraph{Result 2 (the exact price of relaxing confidence)}\label{res:c2} If the demonstration may accept a worse unit
with probability $\delta'>\delta$ when the simulator is wrong---that is, if one accepts a higher consumer's risk---the
real demands can be reduced by at most
\[c=n_0(\varepsilon,\delta)-n_0(\varepsilon,\delta'),\]
and this reduction is attained by a two-tier plan: run the zero-failure test with $n_0(\varepsilon,\delta')$ demands when
a simulator test passes and with $n_0(\varepsilon,\delta)$ otherwise. The plan keeps confidence $1-\delta$ only if a simulator of a worse-than-target unit passes with probability at most
$\pi^\ast$: \nPiStarA, \nPiStarB, \nPiStarC{} and \nPiStarD{} for $\delta'=7.5$, 10, 15 and 20\% at $\varepsilon=\delta=5\%$.
This requirement is itself a link between simulator and reality, so Result~2 is best read as the price of lowering
the confidence, not as a free credit. The frequently
used continuous form $\lfloor\ln(\delta'/\delta)/(-\ln(1-\varepsilon))\rfloor$ is only a lower bound on $c$.
Table~\ref{tab:c2} gives the numbers.

\begin{table}[t]
\centering
\caption{Result 2: exact credit $c$ versus the continuous form $c^\ast$ ($\varepsilon=\delta=5\%$, $n_0=\nNzero$), and the
two-tier plan run with the twin event pool as an independent simulator (worst case $p=\varepsilon$; closed form given the
pass probability, and 100{,}000 Monte-Carlo runs).}
\label{tab:c2}
\begin{tabular}{@{}lcccc@{}}
\toprule
$\delta'$ & 7.5\% & 10\% & 15\% & 20\% \\
\midrule
exact credit $c$ (demands) & \ncExactA & \ncExactB & \ncExactC & \ncExactD \\
saving w.r.t.\ $n_0$ & \ncSaveA & \ncSaveB & \ncSaveC & \ncSaveD \\
continuous form $c^\ast$ & \ncStarA & \ncStarB & \ncStarC & \ncStarD \\
false acceptance, twin pool (closed form) & \nCtwoCFA & \nCtwoCFB & \nCtwoCFC & \nCtwoCFD \\
\quad Monte Carlo & \nCtwoFCA & \nCtwoFCB & \nCtwoFCC & \nCtwoFCD \\
\bottomrule
\end{tabular}
\end{table}

Monte-Carlo runs match the closed form within \nCtwoMaxDevPP{} percentage points on a \nCtwoNgrid-point grid. Used as an
independent simulator (its events resampled independently of the real draws), the twin event pool passes the simulator
test in \nCtwoTwinPass{} of runs, the credit is almost fully realised, and the false-acceptance probability at
$p=\varepsilon$ is $\delta'$ up to the pass probability: the credit is bought exactly with the extra error. An
off-the-shelf synthetic benchmark (BMC~\cite{bmc2012}; \nCtwoBMCN{} events long enough for the operating point) never passes
at $\varepsilon=5\%$ (\nCtwoBMCPass): it has too few synthetic events to bound its own miss probability below $5\%$. At
$\varepsilon=20\%$ it passes in every run and buys the same credit as the twin (supplementary Table~S13), which shows
that passing the simulator test is no evidence of fidelity. Monte-Carlo standard errors are about \nCtwoSE.

\paragraph{Result 3 (paired-twin bound)}\label{res:c3} For any pairing of real and twin events, the real miss
probability is at most the twin miss probability plus the discordance,
\[p\le q+D,\]
with no assumption on the simulator (the bound splits the real misses into those the twin also misses and those it
catches). On a single camera the bound does not save demands: if $q\le q_U$ holds with confidence $1-\delta/2$ and $D\le\varepsilon-q_U$ is
demonstrated by a zero-discordance run at level $\delta/2$, the camera needs
$\lceil\ln(\delta/2)/\ln(1-\varepsilon+q_U)\rceil$ paired events, i.e.\ \nCthreeA, \nCthreeB{} and \nCthreeC{} for
$q_U=0,1,2\%$ at $\varepsilon=\delta=5\%$, all above $n_0=\nNzero$. A twin built with a pessimistic (shorter) event length
remains valid: its miss probability can only grow and its discordance only shrink. On the fixed-camera twin, shortening
every event by 50\% raises the twin miss rate at \OPs{} from \nCfiveQzero{} to \nCfiveQfifty. With a uniform phase the
discordance probability of an event is $|R'\setminus R|/K$, where $R,R'$ are the residues modulo $K$ of the real and twin
hit frames; the worst-phase indicator $\mathbf 1\{R'\setminus R\ne\emptyset\}$ dominates it, so the bounds of Result~4 also
hold for the integer worst-phase counts we report.

\paragraph{Result 4 (transfer across scenes)}\label{res:c4} Let $m$ calibration scenes contribute $N=\sum_sn_s$ paired
events, $X$ of them discordant, and let $J$ scenes show at least one discordant event.
(a)~\emph{Population bound (the transfer guarantee).} If the scenes are drawn independently from a scene population (events within a scene may be dependent), the population-average
discordance of a new camera is at most $\UCP(J,m)$ with confidence $1-\delta$, so its expected miss probability is at most
its expected twin miss probability plus $\UCP(J,m)$. The bound counts scenes, not events; it needs no independence of
events within a scene. Its price is the number of scenes: with $J=0$, $\UCP\le10\%$ needs \nPopScenesTen{} scenes and
$\UCP\le5\%$ needs \nPopScenesFive. With the \nMainScenes{} fixed-camera scenes of this study, $J=0,1,2,3$ give
$\UCP(J,m)=\nUpopJzero$, \nUpopJone, \nUpopJtwo{} and \nUpopJthree{} at level $\delta$ and \nUhalfJzero, \nUhalfJone,
\nUhalfJtwo{} and \nUhalfJthree{} at level $\delta/2$: a single discordant scene costs about two percentage points, and a
bound below $\varepsilon/2=10\%$ tolerates three discordant scenes at level $\delta$ but only two at $\delta/2$. Adding
events to existing scenes does not tighten it; only new scenes do.
(b)~\emph{Fleet bound (calibrated cameras, independent events).} For the calibrated cameras, the event-weighted discordance is at most $\UCP(X,N)$ with confidence
$1-\delta$ (for $\delta\le 1/4$), however heterogeneous the cameras are; hence their event-weighted miss rate is at most
their twin miss rate plus $\UCP(X,N)$. Applying it requires a separate bound on the twin miss rate, with its own share of $\delta$; where real misses exist,
direct Clopper--Pearson on them is at least as informative. The proof compares the Poisson-binomial count with a
binomial of the same mean~\cite{hoeffding1956} and uses a lower bound on the probability that a binomial reaches its mean~\cite{greenberg2014}, applied to the complementary count.
This bound assumes that events are independent given the scenes, so positive dependence within a scene makes
it optimistic, and it says nothing about a new camera.
(c)~\emph{Per-camera limit.} A per-camera statement, ``$p\le q+t$ except on a set of cameras of mass $\gamma$'', follows if
every scene has at least $n_{\min}$ events, with $\gamma=\UCP(J,m)/(1-(1-t)^{n_{\min}})$. Table~\ref{tab:c4b} turns this
into a data requirement: for $\gamma=5\%$ at least \nPopScenesFive{} scenes, each with tens to hundreds of events. No
public benchmark we know offers this, so the per-camera form is reported as a limit.

\begin{table}[t]
\centering
\caption{Data needed for the per-camera form of Result 4(c) with no discordance ($J=0$): minimum number of scenes $m_{\min}$, events per scene $n$ at $m=m_{\min}$ and at $m=2m_{\min}$, total events, and the asymptotic floor $\ln(1/\delta)/(\gamma t)$.}
\label{tab:c4b}
\scriptsize
\setlength{\tabcolsep}{3pt}
\begin{tabular}{@{}rrrrrrr@{}}
\toprule
$t$ & $\gamma$ & $m_{\min}$ & $n$ at $m_{\min}$ & $n$ at $2m_{\min}$ & total at $2m_{\min}$ & floor \\
\midrule
2.5\% & 5\% & 59 & 183 & 28, 3,304 & 2,397 \\
2.5\% & 10\% & 29 & 158 & 28, 1,624 & 1,198 \\
5\% & 5\% & 59 & 91 & 14, 1,652 & 1,198 \\
5\% & 10\% & 29 & 78 & 14, 812 & 599 \\
10\% & 5\% & 59 & 44 & 7, 826 & 599 \\
10\% & 10\% & 29 & 38 & 7, 406 & 300 \\
\bottomrule
\end{tabular}
\end{table}


\paragraph{Demonstration procedure} The procedure evaluated in Section~\ref{sec:results} is: (1)~declare the domain and
pair every real event of $m$ calibration scenes from it with its twin; count discordant scenes $J$; (2)~\emph{domain rule}:
if $\UCP_\delta(J,m)>\varepsilon/2$, withdraw the transfer claim for the whole domain (a design rule for go/no-go, not a
theorem); (3)~for a new camera of the declared domain, run its twin until its own miss probability is bounded, $q_U$, at
level $\delta/2$ (a twin \emph{run} is one rendering of one synthetic event in one sprite variant; runs are treated as
independent, which is optimistic when runs share a trajectory); (4)~accept the camera with no real demand if $q_U+\UCP_{\delta/2}(J,m)\le\varepsilon$, otherwise fall
back to the zero-failure test (Fig.~\ref{fig:scheme}). By Result~4(a) and a union bound, with confidence $1-\delta$ the population-average discordance is at most
$\UCP_{\delta/2}(J,m)$. Among \emph{accepted} cameras the average discordance can be larger, by at most the factor
$1/\PP(\mathrm{accept})$, because acceptance selects cameras with a low twin miss rate; the average miss probability of
accepted cameras is therefore at most $\varepsilon-\UCP_{\delta/2}+\UCP_{\delta/2}/\PP(\mathrm{accept})$. At \OPs{}, where
\nPAccept{} of the static cameras are accepted, the inflation is \nAcceptInflation{} percentage points; out of domain it is
unbounded (Section~\ref{sec:audit}). The acceptance is \emph{not} a per-camera guarantee. In the audit we nevertheless count, per camera, acceptances with $p_s>\varepsilon$ (``false acceptances'') as a
diagnostic of how the average guarantee behaves on individual cameras.

\begin{figure}[t]
\centering
\begin{tikzpicture}[font=\footnotesize, node distance=4mm and 6mm,
  box/.style={draw=black!60, rounded corners=2pt, align=center, inner sep=3pt, minimum height=8mm},
  arr/.style={-{Latex[length=1.8mm]}, draw=black!70}]
\node[box] (real) {real event\\(frames, ground truth)};
\node[box, right=of real] (twin) {paired twin\\(background plate + sprite)};
\node[box, below=of real] (M) {$M$: real miss\\at phase $\varphi$};
\node[box, below=of twin] (Mp) {$M'$: twin miss\\at the same $\varphi$};
\node[box, right=8mm of twin, yshift=-6mm] (qd) {$q$, $D=\PP(M{=}1,M'{=}0)$\\per scene};
\node[box, below=of qd, yshift=-2mm] (cert) {fleet bound $\UCP(X,N)$\\population bound $\UCP(J,m)$};
\node[box, below=9mm of M, xshift=20mm, minimum width=60mm] (dec) {new camera: accept iff $q_{\mathrm{new}}+\UCP\le\varepsilon$};
\draw[arr] (real) -- (twin);
\draw[arr] (real) -- (M);
\draw[arr] (twin) -- (Mp);
\draw[arr] (M) -- (Mp);
\draw[arr] (Mp.east) -- (qd.west);
\draw[arr] (qd) -- (cert);
\draw[arr] (cert.south) |- (dec.east);
\end{tikzpicture}
\caption{The demonstration with a paired twin. Each real event is paired with a synthetic twin that replaces the object and
keeps the scene; both are evaluated by the same detector at the same sampling phase. Discordances calibrated across scenes
bound the miss probability of a new camera through $p\le q+D$.}
\label{fig:scheme}
\end{figure}
